\documentclass[11pt,a4paper]{article}

% ====== PACKAGES ======
\usepackage[utf8]{inputenc}
\usepackage[T1]{fontenc}
\usepackage[english]{babel}
\usepackage{geometry}
\geometry{top=2.5cm, bottom=2.5cm, left=2.5cm, right=2.5cm}

\usepackage{amsmath, amsfonts, amssymb}
\usepackage{graphicx}
\usepackage{hyperref}
\usepackage{booktabs}
\usepackage{caption}
\usepackage{subcaption}
\usepackage{float}
\usepackage{xcolor}
\usepackage{listings}
\usepackage{tcolorbox}
\usepackage{algorithm}
\usepackage{algpseudocode}
\usepackage{enumitem}
\usepackage{setspace}
\usepackage{titlesec}
\newcommand{\sectionbreak}{\clearpage}
\usepackage{listings}
\renewcommand*{\lstlistingname}{Code Excerpt}


\onehalfspacing

% Link configuration
\hypersetup{
    colorlinks=true,
    linkcolor=blue!60!black,
    filecolor=magenta,      
    urlcolor=cyan,
    pdftitle={DSAI Project Report},
}

% Code block configuration
\definecolor{codegreen}{rgb}{0,0.6,0}
\definecolor{codegray}{rgb}{0.5,0.5,0.5}
\definecolor{codepurple}{rgb}{0.58,0,0.82}
\definecolor{backcolour}{rgb}{0.95,0.95,0.92}

\lstdefinestyle{mystyle}{
    backgroundcolor=\color{backcolour},   
    commentstyle=\color{codegreen},
    keywordstyle=\color{magenta},
    numberstyle=\tiny\color{codegray},
    stringstyle=\color{codepurple},
    basicstyle=\ttfamily\footnotesize,
    breakatwhitespace=false,         
    breaklines=true,                 
    captionpos=b,                    
    keepspaces=true,                 
    numbers=left,                    
    numbersep=5pt,                  
    showspaces=false,                
    showstringspaces=false,
    showtabs=false,                  
    tabsize=2
}
\lstset{style=mystyle}

\title{\Huge \textbf{DSAI Project} \\ \vspace{0.5cm}
\Large Debate Analysis: Convincing Argument Detection, AI-Generated Text Detection, and Debate Generation with Fine-tuned LLMs}

\author{
\begin{tabular}{c}
\textbf{Authors:}\\
SMATI Wassim \\
DAHASSE Yanis \\
CASTEL Nicolas \\
JIN Tristan \\[0.8cm]
\textbf{Supervisors:}\\
LABEAU Matthieu \\
FIHEY Pierre
\end{tabular}
}

\date{}

\begin{document}

\maketitle



\newpage
\tableofcontents
\newpage

% ==========================================
% SECTION 1: INTRODUCTION
% ==========================================
\newpage
\section{Introduction}
\subsection{Project Background}
This project was not among the topics initially proposed on Moodle.
Our group wanted to explore and deepen our knowledge of NLP following the
Machine Learning for Text Mining course taught by Matthieu Labeau.
Since project selection was competitive, we took the initiative of contacting
him directly to propose an open-ended project.\\
Our early ideas focused on empathy assessment in health-related discourse,
pedagogical analysis of videos and articles, or greenwashing detection in
ambiguous environmental communications. After discussing with our supervisor,
we converged instead on the computational analysis of online debates.\\
This project uses the \textbf{ChangeMyView} (CMV) subreddit as its main data source,
a platform where users share their opinions and award a $\Delta$ symbol to
arguments that successfully changed their mind. This validation mechanism
provides a natural label that is rare in web data, making CMV an ideal corpus
for studying persuasion. From this framework, we structured the project into
three axes, presented below.

\subsection{Project Structure}
The project follows a logical progression: understanding persuasion,
detecting synthetic content, and then applying both components to arguments
generated by fine-tuned LLMs.

The \textbf{first axis} focuses on the study of argument persuasiveness. Using the
\textbf{Winning Argument Corpus} (WAC), a filtered and cleaned version of CMV,
we aim to build a classifier capable of distinguishing a persuasive argument
from a non-persuasive one. We explore several approaches: classic vector
encoders (TF-IDF, Word2Vec, RoBERTa) followed by an SVM, a fine-tuned
RoBERTa, and linguistic features inspired by the work of Tan et al.~\cite{tan2016winning}
capturing both the intrinsic style of the argument and its relationship with
the original post. This classifier forms the core building block of the project,
and will be reused in the third axis.

A natural complementary question arises: if language models can produce
convincing arguments, can we tell them apart from human-written text? This
is the challenge of the \textbf{second axis}, which addresses AI-generated text
detection. Relying on three corpora of increasing difficulty , GriD, HC3
and M4GT-Bench , we evaluate several families of methods, from
TF-IDF~+~SVM to full RoBERTa fine-tuning. The goal is to measure not only
the raw performance of each approach, but also its ability to generalise
across different generators and domains. This model will serve as an
authenticity detector in the third axis.

The goal of the \textbf{third axis} is to fine-tune an open-source LLM (GPT-2,
Qwen) on WAC data, then steer its generation to simultaneously maximise
the persuasion score from Axis 1 and the authenticity estimated by Axis 2.
Two strategies are explored: \textit{prompt engineering} guided by the
classifier coefficients, and \textit{Best-of-N} sampling, where the best
argument is chosen from several generated candidates. Due to hardware
constraints, fine-tuning a 7-billion-parameter model on a consumer GPU
led us to implement QLoRA to make this process tractable.

% ==========================================
% SECTION 3: AXIS 1
\newpage

\section{Axis 1: Modelling and Predicting Persuasion}

In this section, we explore several ways to assess the level of persuasiveness
of an argument. Our goal at the end of this axis is to obtain the best possible
classification model, capable of assigning a score and classifying an argument
as persuasive or not. To this end, we use the Winning Argument Corpus (WAC),
along with encoder-based methods such as W2V, TF-IDF or RoBERTa followed by
an SVM or a fine-tuned linear layer, as well as less conventional approaches
such as the pairwise method.



\subsection{Winning Argument Corpus (WAC)}
\subsubsection{Dataset Presentation}
The \emph{Winning Argument Corpus} is built from the \textbf{r/ChangeMyView}
subreddit, where an author presents an opinion (the \emph{original post}, OP)
and other users reply to try to change their mind. If the OP considers that
a reply has changed their position, they award it a \textbf{$\Delta$ (delta)}
symbol. Each discussion thus provides pairs of ``winning'' / ``losing'' replies
to the same post, used to study what makes an argument persuasive. Below is
an excerpt from a discussion taken from the dataset:

\vspace{0.3cm}

\begin{figure}[H]
\begin{lstlisting}[basicstyle=\ttfamily\scriptsize, breaklines=true, frame=single]
{"id": "t3_2cm3oh", "root": "t3_2cm3oh", "user": "nintendosixtyfooour",
 "reply-to": null, "text": "I go to the movies a lot and it seems to happen...",
 "meta": {"success": null, "pair_ids": [], "score": null, ...}}

{"id": "t1_cjgsk1z", "root": "t3_2cm3oh", "user": "KevinWestern",
 "reply-to": "t3_2cm3oh", "text": "Haven't you ever been near a sports fan...",
 "meta": {"success": 1, "pair_ids": ["p_779"], "score": 757, ...}}

{"id": "t1_cjgv6jc", "root": "t3_2cm3oh", "user": "Notsuchameangirl",
 "reply-to": "t3_2cm3oh", "text": "Clapping at the end of a movie you enjoyed...",
 "meta": {"success": 0, "pair_ids": ["p_779"], "score": 9, ...}}

{"id": "t1_cjh1biu", "root": "t3_2cm3oh", "user": "TomasTTEngin",
 "reply-to": "t3_2cm3oh", "text": "It's an effective way to remove popcorn residue...",
 "meta": {"success": null, "pair_ids": [], "score": 26, ...}}
\end{lstlisting}
\caption{Excerpt from \texttt{utterances.jsonl} (a WAC file)}
\label{fig:wac-example}
\end{figure}

In \texttt{utterances.jsonl}, the \texttt{meta.success} field indicates the
status of each comment relative to the OP: \texttt{1} marks a comment that
received a delta (\textbf{persuasive}) and was included in an argument pair,
and \texttt{0} marks the other member of that pair, which did not receive a
delta (\textbf{non-persuasive}) and serves as a comparison point. The
\texttt{null} label is by far the most frequent, covering any comment not
retained in a pair (either because it is not a genuine persuasion attempt or
because it cannot be meaningfully compared to the pair's argument). The link
between the two members of a pair is given by \texttt{meta.pair\_ids}. As for
\texttt{meta.score}, it represents the number of upvotes a comment received;
we do not use it here.

\subsubsection{Simplification of Axis 1 \& Dataset Cleaning}

The original goal of Axis 1 was to detect the $\Delta$ among all comments in
a post. After considering several approaches, we decided to significantly
simplify the problem by framing it as a binary classification task: for each
original post, we only keep the \texttt{0/1} argument pairs, discarding all
comments labelled \texttt{null}. We can then build a model capable of
distinguishing a winning comment from a losing one.\\

This simplification is nevertheless quite rough and discards a significant
amount of information. The entire tree structure of the debate disappears:
by removing \texttt{null}-labelled comments, we lose the chains of
argumentation and rebuttal that sometimes lead to a $\Delta$, as well as the
temporal position of each comment in the discussion. The score (upvote count)
of each comment is also ignored, even though it could provide complementary
information, albeit distinct from the OP's persuasion.

\vspace{0.3cm}

\begin{figure}[H]
\begin{lstlisting}[basicstyle=\ttfamily\scriptsize, breaklines=true, frame=single]
id,root,op_text,text,success,pair_id
t1_cjgsk1z,t3_2cm3oh,"I go to the movies a lot and it seems to happen...",
  "Haven't you ever been near a sports fan...",1.0,p_779
t1_cjgv6jc,t3_2cm3oh,"I go to the movies a lot and it seems to happen...",
  "Clapping at the end of a movie you enjoyed...",0.0,p_779
\end{lstlisting}
\caption{Excerpt from \texttt{WAC.csv} , the two rows forming pair \texttt{p\_779}}
\label{fig:wac-simplifie}
\end{figure}

Here is an excerpt from our filtered dataset in CSV format, with 6 columns:
\texttt{id} (comment identifier), \texttt{root} (identifier of the original
post being replied to), \texttt{op\_text} (the original post), \texttt{text}
(the comment), \texttt{success} (winning or losing label) and \texttt{pair\_id}
(identifier linking the two comments of the same pair). In total, we obtain
3866 winning/losing argument pairs, giving a dataset of \textbf{7732} perfectly
balanced arguments, which justifies our use of standard Accuracy to evaluate
our models going forward.


\subsection{Approach Using Classic Vector Representations}
\subsubsection{Encoders : TF-IDF, Word2Vec \& RoBERTa}

Our first approach consists of vectorising comments using different classic
encoders, then training an SVM on the resulting vectors. The representations
produced by these encoders are compared to each other, in order to assess
how much the semantic richness of an encoder influences the SVM's ability to
distinguish a persuasive comment from a non-persuasive one. The three
encoders tested are:

\begin{itemize}
    \item \textbf{TF-IDF}: A representation based on word frequency, which
    captures discriminating vocabulary without taking word order or context
    into account.
    \item \textbf{Word2Vec}: A model that learns a dense vector for each word
    from the words that frequently surround it in the text, so that words with
    similar meanings obtain close vectors in the vector space.
    \item \textbf{RoBERTa}: A pre-trained Transformer model that produces, for
    each word, a context-dependent vector, which allows it to capture nuances
    of meaning that more static approaches like Word2Vec cannot distinguish.
\end{itemize}

For the last two encoders, the final vector for a comment is obtained by
averaging the representations of each word in the comment into a single vector.

\subsubsection{Results and Discussion}

We run tests according to two strategies. The first consists of vectorising
only the comment, completely ignoring the original post (OP): this evaluates
the intrinsic quality of the argument independently of the debate context.
The second consists of concatenating the OP and the comment before vectorising
the whole thing, which adds the debate context to the comment, a crucial
element when assessing persuasiveness, since the same argument may be more or
less relevant depending on the OP's initial position. In both cases, training
is performed on 80\% of the dataset, with the remaining 20\% reserved for
testing. The results are as follows:
\vspace{0.2cm}
\begin{table}[h]
\centering
\begin{tabular}{lcc}
\toprule
\textbf{Model} & \textbf{Text only} & \textbf{Text + OP} \\
\midrule
Word2Vec + SVM  & 0.56 & 0.49 \\
TF-IDF + SVM    & 0.51 & 0.32 \\
RoBERTa + SVM   & 0.54 & 0.35 \\
\bottomrule
\end{tabular}
\caption{SVM Accuracy by encoder and input context}
\label{tab:axe1_results}
\end{table}

All results remain close to randomness (0.50) regardless of the encoder, which 
confirms the inherent difficulty of the task: distinguishing a persuasive 
argument from a non-persuasive one purely on the basis of its textual content 
is a fundamentally hard problem, even for rich semantic representations.

%%%%% A MODIF + SUR LES RESULTATS DE CHAQUE ENCODERS

In the text only setting, Word2Vec achieves the best accuracy 
(0.56), marginally ahead of RoBERTa (0.54) and clearly ahead of TF-IDF 
(0.51). This is somewhat counterintuitive: RoBERTa's contextual 
representations are expected to be richer, but averaging over all tokens 
collapses them into a document vector that loses much of that contextual 
information. Word2Vec's averaged embeddings suffer from the same 
aggregation loss, but their denser geometry appears more amenable to 
linear SVM separation on a small dataset.

The most striking finding is the accuracy drop when the OP is 
concatenated to the comment. For TF-IDF, accuracy collapses from 0.51 to 
0.32, well below chance with a similarly effect for RoBERTa 
(0.54 to 0.35). This dilution effect is explained by the pairwise 
structure of the dataset: since both comments in a pair respond to the same 
OP, concatenating it artificially brings their vectors closer together, 
cancelling the discriminative signal. These results show that naive context 
injection via concatenation is counter-productive, and motivate the more 
principled approaches explored in Sections~2.3 and~2.5.


\subsection{Intrinsic Features and Interplay}
\subsubsection{Paper's features}

The results obtained show that vectorising the comment alone or naively
concatenating the OP and the comment does not fully exploit the relationship
between the two texts. We therefore turn to the work of Tan et
al.~\cite{tan2016winning} on the same \textbf{ChangeMyView} corpus, which
proposes around forty linguistic features organised into two families.\\

The first family, called \textbf{interplay}, captures the relationship between
the comment and the original post: for each pair (comment $A$, original post
$O$), they compute the number of shared words $\boldsymbol{|A \cap O|}$,
the reply-side fraction $\boldsymbol{\frac{|A \cap O|}{|A|}}$, the OP-side
fraction $\boldsymbol{\frac{|A \cap O|}{|O|}}$, and the Jaccard similarity
$\boldsymbol{\frac{|A \cap O|}{|A \cup O|}}$, applied to stopwords, content
words and the full vocabulary. The key finding is that persuasive arguments
are less similar to the OP in terms of content words, while remaining more
similar in terms of stopwords. They thus bring new perspectives while
stylistically aligning with their interlocutor.\\

The second family groups \textbf{intrinsic or stylistic} descriptors of the comment:
length (\#words, \#sentences, \#paragraphs), pronoun usage (first-person
singular and plural, second-person), presence of hyperlinks, example markers,
hedges, and Markdown formatting (italics, bold, lists).

\subsubsection{Results and Discussion}


\begin{table}[h]
\centering
\begin{tabular}{lcccc}
\toprule
\textbf{Class} & \textbf{Precision} & \textbf{Recall} & \textbf{F1-score} & \textbf{Support} \\
\midrule
0 (Losing) & 0.58 & 0.70 & 0.63 & 774 \\
1 (Winning)     & 0.62 & 0.50 & 0.54 & 773 \\
\bottomrule
\end{tabular}
\caption{Performance of Paper's Features + SVM; \textbf{Accuracy = 0.59}}
\label{tab:features_svm}
\end{table}

The linguistic features proposed by Tan et al.\ achieve an accuracy of 0.59, 
a clear improvement over all encoder-based approaches from Section~2.2. 
This result is notable given the compactness of the feature vector. A few hand-crafted features outperform dense pre-trained representations, 
suggesting that persuasiveness in this dataset is better captured by structured rhetorical signals than by raw content.

The class-level breakdown reveals an asymmetry: the model recalls losing arguments more reliably (recall 0.70) than winning ones (recall 0.50). The classifier is thus more confident at identifying what makes an argument fail than what makes it succeed. Interplay and stylistic markers can flag obvious weaknesses, but the positive signals of persuasion remain subtler and harder to see with a fixed feature set. The SHAP analysis in Section~2.3.3 investigates which features actually drive these decisions.


\subsubsection{Feature Analysis with SHAP}

To interpret the SVM model, we use the SHAP method. At the local level, it
quantifies the impact of each variable on a prediction: a positive value
pushes towards the ``persuasive'' class (1), and vice versa. At the global
level, aggregating these values over the test set gives the overall feature
importance ranking.

\begin{figure}[h]
    \centering
    \includegraphics[width=0.5\linewidth]{shap_analysis.png}
    \caption{Feature importance representation}
    \label{fig:placeholder}
\end{figure}

\textbf{Interplay} features dominate by a wide margin, which confirms the
findings of Tan et al. An apparent contradiction is worth noting:
a high \texttt{op\_frac\_all} is associated with positive SHAP values, while
a high \texttt{jaccard\_all} generates negative SHAP values. This is explained
by normalisation: \texttt{op\_frac\_all} $= \frac{|A \cap O|}{|O|}$ measures
OP coverage, whereas Jaccard $= \frac{|A \cap O|}{|A \cup O|}$ penalises
longer comments. A persuasive argument therefore covers a large part of the
OP's vocabulary while introducing a lot of new vocabulary , low Jaccard,
high \texttt{op\_frac} , which remains consistent with Tan et al.
A high \texttt{reply\_frac\_all} generates negative SHAP values: paraphrasing
the OP without adding anything new hurts persuasiveness. Among intrinsic
features, \texttt{frac\_links} and \texttt{n\_sentences} contribute positively,
while formatting (\texttt{n\_italics}) and lexical diversity (\texttt{word\_entropy})
have a marginal impact.


\subsection{The Pairwise Method}

This section complements our main approach. We focus here on the pairwise
method, as implemented in Tan et al. Rather than classifying each argument
independently, we feed the classifier directly with winning/losing argument
pairs, and the task consists in identifying which of the two is persuasive.
This reformulation simplifies the problem, since both arguments respond to the
same OP and are therefore directly comparable, which leads us to expect better
performance.

In practice, we encode each argument in the pair with our four vectorisation
methods (TF-IDF, Word2Vec, RoBERTa and Tan et al.'s features), then subtract
the two vectors to obtain a single vector representing the contrast between
the two arguments. This vector is then fed into a Logistic Regression
classifier to predict which argument is more persuasive.

\begin{table}[h]
\centering
\begin{tabular}{lcc}
\toprule
\textbf{Encoder} & \textbf{Classic Approach (SVM)} & \textbf{Pairwise Method (LogReg)} \\
\midrule
TF-IDF           & 0.51 & 0.58 \\
W2V              & 0.56 & 0.59 \\
RoBERTa          & 0.54 & 0.51 \\
Paper's Features & 0.59 & 0.63 \\
\bottomrule
\end{tabular}
\caption{Accuracy comparison between SVM and Pairwise approaches}
\label{tab:encoder_results}
\end{table}

The pairwise reformulation yields consistent gains across all encoders. The improvement is most pronounced for TF-IDF (0.51 to 0.58) 
and Paper's Features (0.59 to 0.63), while RoBERTa is the only encoder 
that does not benefit from the reformulation (0.54 to 0.51). This exception 
is likely due to the vector subtraction step: RoBERTa's contextual 
embeddings are not designed to be compared by difference. Overall, these results confirm that explicitly comparing argument pairs,rather than classifying each 
independently is a more natural framing for the task, consistent with 
the pairwise structure of the WAC dataset itself.

\subsection{Fine-tuned RoBERTa}

The results obtained so far have highlighted a clear limitation: vectorising
arguments, even with a powerful encoder like RoBERTa, is not sufficient to
capture the dynamic relationship between an argument and the original post.
We therefore decided to move to a full fine-tuning of RoBERTa, adapting all
model weights to our pairwise classification task.

\subsubsection{The Multiple Choice Architecture}

We keep the pairwise approach that showed its value in Section 2.4, but
implementing it with a Transformer model raises a concrete problem: the strict
512-token limit. Naively concatenating the OP and both arguments would force
us to truncate the texts heavily, losing crucial information such as the
conclusion of the argument , often where persuasion is decided.

We therefore reformulated the task using the \texttt{RobertaForMultipleChoice}
architecture. Concretely, the model separately evaluates the pair
\texttt{(OP, Argument 0)}, then the pair
\texttt{(OP, Argument 1)}. A final linear layer then compares
the two representations to predict the winner. This architecture allows us to
allocate 512 tokens to each pair, giving twice the budget of a naive
concatenation of the three texts, while letting the model capture the
relationship between the OP and each argument via the Transformer's internal
cross-attention.

\subsubsection{Data Preparation and Bias Handling}

To train this model correctly, we had to adapt our dataset and take several
precautions that seemed important to us.

The first concerns \textbf{position bias}. We randomly shuffled the position
of the winning argument for each pair. Without this precaution, the network
would quickly learn to systematically predict the same position to minimise
its error, without ever reading the text.

The second relates to \textbf{data leakage}. The train/validation/test split
was performed by grouping by post identifier. This ensures that a test argument
always responds to an OP that the model has never seen during training.


\subsubsection{Results and Interpretation}

The final evaluation was conducted on a test set of 387 pairs never seen
during training. The detailed results are presented in Table~\ref{tab:roberta_results}.

\begin{table}[h]
\centering
\begin{tabular}{lcccc}
\hline
\textbf{Class} & \textbf{Precision} & \textbf{Recall} & \textbf{F1-score} & \textbf{Support} \\
\hline
0 (Option 0 wins) & 0{,}69 & 0{,}76 & 0{,}72 & 187 \\
1 (Option 1 wins) & 0{,}75 & 0{,}69 & 0{,}72 & 200 \\
\hline
\end{tabular}
\caption{Performance of the fine-tuned RoBERTa Cross-Encoder trained on 4 epoch; \textbf{Accuracy = 0{,}72}}
\label{tab:roberta_results}
\end{table}

The fine-tuned RoBERTa Cross-Encoder achieves an accuracy of 72\%, a substantial improvement over all previous approaches. Both classes obtain a balanced F1-score of 0.72, confirming that the model does not exhibit the recall asymmetry observed with hand-crafted features in Section~2.3.

This result validates the Cross-Encoder architecture: by encoding each 
(OP, argument) pair separately within a 512-token budget, the model can 
attend to the relationship between the post and each candidate without 
suffering from the dilution effect identified in Section~2.2. Full 
fine-tuning thus allows RoBERTa to learn persuasion-specific patterns directly from raw text, effectively replacing the features 
 of Tan et al's paper\ while surpassing it by a clear margin 
(0.72 vs.\ 0.63). This model constitutes our strongest 
scoring function and will be used as such in the Axis~3 generation 
pipeline. 
\vspace{1cm}

\begin{table}[h]
\centering
\begin{tabular}{lcc}
\hline
\textbf{Model} & \textbf{Approach} & \textbf{Accuracy} \\
\hline
TF-IDF + SVM          & Classic    & 0{,}51 \\
Word2Vec + SVM        & Classic    & 0{,}56 \\
RoBERTa + SVM         & Classic    & 0{,}54 \\
Paper's Features + SVM & Classic  & 0{,}59 \\
\hline
TF-IDF                & Pairwise   & 0{,}58 \\
Word2Vec              & Pairwise   & 0{,}59 \\
Paper's Features      & Pairwise   & 0{,}63 \\
\hline
\textbf{Fine-tuned RoBERTa} & \textbf{Cross-Encoder} & \textbf{0{,}72} \\
\hline
\end{tabular}
\caption{Summary of all Axis 1 approaches on the WAC corpus}
\label{tab:axe1_recap}
\end{table}


\newpage

\section{Axis 2: Automated Detection (Human vs. AI)}

Axis 2 addresses a binary safety classification problem: determining whether
a text was written by a human or generated by an artificial intelligence model.
In the overall project architecture, this module acts as an authenticity filter:
it assesses whether the arguments produced or analysed retain a sufficiently
human-like signature.

Rather than presenting each experiment separately, we group the results here
according to a comparative logic: the same encoder families, the same linear
classifiers where possible, but applied to three datasets of increasing difficulty.
This organisation makes the strengths and limitations of each approach easier
to see.

\subsection{Data Corpora: HC3, GriD and M4GT}
The evaluation relies on three complementary corpora, chosen to test both raw
performance and the generalisation ability of the detectors:
\begin{itemize}
    \item \textbf{HC3 (Human ChatGPT Comparison Corpus):} A corpus built around
    pairs of human and ChatGPT responses. It provides a useful framework for
    measuring the separability between human style and the style of a
    conversational model.
    \item \textbf{GriD:} A dataset used as an intermediate benchmark, allowing
    comparison of lexical, semantic and contextual representations on a smaller
    number of examples.
    \item \textbf{M4GT-Bench:} A more heterogeneous corpus, integrating multiple
    generators and multiple domains. It is therefore more representative of an
    \textit{Out-Of-Distribution} scenario, where the detector must not only
    recognise a specific model, but generalise to varied generation styles.
\end{itemize}

\subsection{Methods Evaluated}
We compared four families of approaches, ranging from simple statistical
representations to full Transformer fine-tuning:

\subsubsection{TF-IDF + SVM}
The TF-IDF method transforms each text into a high-dimensional sparse vector,
where each component represents the relative importance of a term in the
document and in the corpus. This representation is then fed to a linear SVM.
Its main advantage is its simplicity, speed and robustness on large text
corpora. It captures differences in vocabulary, phrasing and lexical frequency
between human and generated texts very well.

\subsubsection{Word2Vec + SVM}
Word2Vec projects words into a dense space where geometric proximities reflect
semantic similarities. To represent a full document, word vectors are aggregated
and then passed to a linear SVM. This method goes beyond simple lexical
frequency, but loses a significant part of syntactic information and word order.

\subsubsection{SentenceTransformer / RoBERTa + SVM}
Models such as SentenceTransformer or RoBERTa produce contextual embeddings:
the representation of a word or sentence depends on the global context. These
vectors are then used as input to an SVM. This approach aims to capture finer
signals than TF-IDF or Word2Vec, but its performance depends heavily on the
quality of the embedding and its adaptation to the domain.

\subsubsection{RoBERTa Fine-tuning}
Finally, we tested a full fine-tuning of \texttt{roberta-base} on GriD using
\texttt{AutoModelForSequenceClassification}. Unlike RoBERTa + SVM, the model
is not only used as a feature extractor: its weights are adjusted to directly
optimise human/AI classification. This approach is more costly, but can capture
contextual and stylistic signals much more specific to the dataset.

\subsection{Comparative Experimental Results}
The table below summarises all the results extracted from the classification
reports. To keep the overview compact, we report the overall accuracy,
per-class F1-scores and total support. Detailed precision and recall values
are then discussed in the per-dataset analysis.

\begin{table}[H]
    \centering
    \renewcommand{\arraystretch}{1.3}
    \caption{Global comparison of AI detection methods on GriD, HC3 and M4GT.}
    \label{tab:axe2-comparatif-global}
    \resizebox{\textwidth}{!}{%
    \begin{tabular}{llcccccl}
        \toprule
        Dataset & Method & Accuracy & F1 class 0 & F1 class 1 & Support 0 & Support 1 & Main observation \\
        \midrule
        GriD & TF-IDF + SVM & 0{,}93 & 0{,}96 & 0{,}81 & 1053 & 250 & Very good lexical baseline \\
        GriD & Word2Vec + SVM & 0{,}91 & 0{,}94 & 0{,}73 & 1053 & 250 & Lower recall on class 1 \\
        GriD & RoBERTa fine-tuning & 0{,}99 & 1{,}00 & 0{,}98 & 1076 & 227 & Best method on GriD \\
        \midrule
        HC3 & TF-IDF + SVM & 0{,}98 & 0{,}99 & 0{,}97 & 10268 & 3332 & Very strong lexical separation \\
        HC3 & Word2Vec + SVM & 1{,}00 & 1{,}00 & 1{,}00 & 10268 & 3332 & Near-perfect separation \\
        HC3 & RoBERTa + SVM & 0{,}95 & 0{,}97 & 0{,}91 & 10268 & 3332 & Good score despite mixed t-SNE \\
        \midrule
        M4GT & TF-IDF + SVM & 0{,}88 & 0{,}85 & 0{,}89 & 1667 & 2333 & Best method on M4GT \\
        M4GT & Word2Vec + SVM & 0{,}76 & 0{,}68 & 0{,}81 & 1667 & 2333 & Harder generalisation \\
        M4GT & RoBERTa + SVM & 0{,}83 & 0{,}80 & 0{,}85 & 1667 & 2333 & Trade-off between TF-IDF and W2V \\
        \bottomrule
    \end{tabular}%
    }
\end{table}

\subsection{t-SNE Visualisation on HC3: Methodological Choice}

To complement the metrics, we projected the high-dimensional HC3 representations into two dimensions. We specifically chose \textbf{t-SNE} (t-Distributed Stochastic Neighbor Embedding) over standard PCA (Principal Component Analysis). PCA is a linear projection that attempts to preserve global variance, which often results in heavily overlapping clusters when dealing with complex, non-linear text embeddings. Conversely, t-SNE excels at preserving local neighbourhoods and non-linear relationships, making it the ideal tool to visually reveal whether distinct semantic clusters (Human vs AI) exist natively within the latent space of our encoders.

These visualisations give useful intuition about the structure of the vector space learned by each encoder...

\begin{figure}[H]
    \centering
    \begin{subfigure}{0.30\textwidth}
        \centering
        \includegraphics[width=\textwidth]{figures/hc3_tfidf_tsne.png}
        \caption{TF-IDF}
        \label{fig:hc3-tfidf-tsne}
    \end{subfigure}
    \hfill
    \begin{subfigure}{0.30\textwidth}
        \centering
        \includegraphics[width=\textwidth]{figures/hc3_w2v_tsne.png}
        \caption{Word2Vec}
        \label{fig:hc3-w2v-tsne}
    \end{subfigure}
    \hfill
    \begin{subfigure}{0.30\textwidth}
        \centering
        \includegraphics[width=\textwidth]{figures/hc3_roberta_tsne.png}
        \caption{RoBERTa}
        \label{fig:hc3-roberta-tsne}
    \end{subfigure}
    \caption{t-SNE projections of TF-IDF, Word2Vec and RoBERTa representations on HC3.}
    \label{fig:hc3-tsne-comparaison}
\end{figure}

\noindent The TF-IDF projection shows a partial separation: the classes are not fully
disjoint, but some regions are predominantly populated by one category.
Word2Vec gives a much cleaner separation, consistent with its accuracy rounded
to 100\%. RoBERTa produces a more mixed space: the classes overlap more in the
2D projection, which partly explains its lower score on HC3 when used as a
simple embedding extractor followed by an SVM.

\subsection{Per-Dataset Analysis}

\subsubsection{GriD: advantage to supervised fine-tuning}
On GriD, the TF-IDF + SVM baseline already achieves 93\% accuracy, with an
F1-score of 0.96 for class 0 and 0.81 for class 1. The MiniLM + SVM model
drops slightly to 91\% accuracy, mainly due to lower recall on the minority
class. RoBERTa fine-tuning clearly dominates with 99\% accuracy and F1-scores
of 1.00 and 0.98. This shows that, on a controlled-size corpus, supervised
adjustment of a Transformer can capture highly discriminating contextual
signatures.

\subsubsection{HC3: Strong separability and the ChatGPT signature}
HC3 is the dataset on which scores are highest. TF-IDF + SVM achieves 98\%
accuracy, while Word2Vec + SVM achieves an accuracy rounded to 100\%. 
While these results seem exceptional, a critical analysis of the corpus explains this phenomenon. HC3 was built using early versions of ChatGPT (GPT-3.5). This model exhibits a highly formatted, recognisable, and homogeneous "robotic" style (frequent use of bullet points, polite disclaimers, and lack of emotional variability). 

Consequently, the lexical distribution of ChatGPT in HC3 is extremely distinct from human internet users. Even a basic lexical method like TF-IDF easily picks up on these repetitive patterns. This near-perfect score is therefore more indicative of the dataset's simplicity and the stereotypical nature of early ChatGPT, rather than the intrinsic robustness of our detectors.

\subsubsection{M4GT: The Out-Of-Distribution generalisation challenge}
M4GT is the most challenging corpus, and the performance drop is significant:
TF-IDF + SVM drops to 88\% accuracy, RoBERTa + SVM to 83\%, and Word2Vec + SVM falls to 76\%. 

We can deduce several critical insights from this drop. Unlike HC3, M4GT is a multi-generator corpus (Llama, Mistral, Cohere, etc.) covering multiple domains. The "AI signature" is no longer a single homogeneous style. Word2Vec struggles severely because modern LLMs have learned to use a vocabulary that perfectly mimics human semantic distribution. TF-IDF resists slightly better because it captures broader n-gram frequencies (syntactic habits rather than pure semantics). M4GT highlights the core vulnerability of AI detection: detectors that overfit to a specific model's stylistic quirks (like ChatGPT in HC3) fail to generalise to the broader, out-of-distribution (OOD) landscape of modern generative AI.

\subsection{Summary and Pipeline Choice}
These experiments show that AI detection does not come down to choosing the
most complex model. TF-IDF, despite its simplicity, remains extremely
competitive, especially on M4GT where it outperforms dense embeddings. Word2Vec
works remarkably well on HC3, but becomes less reliable on M4GT. Fine-tuned
RoBERTa is excellent on GriD, but RoBERTa used as a simple encoder followed
by an SVM does not systematically dominate.

For the project's full pipeline, this analysis justifies the use of RoBERTa
as an authenticity judge when a fine-tuned model is available, while keeping
TF-IDF as a robust, low-cost baseline. The most important result is methodological:
strong performance on a homogeneous dataset like HC3 does not guarantee reliable
generalisation on a more open corpus like M4GT.

An important theoretical conclusion emerges regarding the use of deep contextual models. Using RoBERTa as a "frozen" feature extractor (RoBERTa + SVM) yielded mixed results (e.g., 83\% on M4GT, behind TF-IDF). We deduce that pre-trained RoBERTa embeddings are fundamentally optimised for semantic similarity (Masked Language Modeling), not stylometry or authenticity. A human text and an AI text expressing the exact same idea will be projected very close to each other in a frozen semantic space, confusing the linear SVM.

However, full \textbf{Fine-Tuning} of RoBERTa (which achieved 99\% on GriD) completely reorganises this embedding space. By adjusting the internal attention weights for a Sequence Classification task, the model learns to break semantic overlaps and focus on structural authenticity markers. This justifies our architectural decision to always favour a fine-tuned Transformer over a static encoder + SVM pipeline when computational resources permit.

% ==========================================
% SECTION 5: AXIS 3
% ==========================================
\newpage

\section{Axis 3: Strategic Generation and LLM Fine-tuning}

Axis 3 represents the convergence of our previous efforts. The goal is twofold:
generate arguments capable of shifting human opinion (Axis 1), while
maintaining a stylistic authenticity that is undetectable by our own detection
algorithms (Axis 2).

\subsection{Causal Fine-Tuning: Architectures and Constraints}
Generation requires an autoregressive model (\textit{Decoder-only Transformer})
trained under the causal language modelling objective (\textit{Causal Language
Modeling}), which aims to maximise the probability of the next token given the
previous tokens:
\begin{equation}
    \mathcal{L}_{CLM} = -\sum_{t=1}^T \log P(w_t | w_{<t}, \theta)
\end{equation}

Two models were experimented with: \textbf{GPT-2} (124M parameters) for rapid
prototyping iterations, and \textbf{Mistral-7B-Instruct-v0.2} (7 billion
parameters) for state-of-the-art inference. Fine-tuning Mistral required
overcoming the dreaded ``CUDA Out of Memory'' (OOM) error on the RTX 3090
(24 GB VRAM) of our server \texttt{nodemm01}.

\subsubsection{Optimisation via QLoRA}
Full parameter adaptation of a 7B model with the AdamW optimiser requires more
than 100 GB of VRAM. To make this process feasible, the \textbf{QLoRA}
(\textit{Quantized Low-Rank Adaptation}) method was rigorously implemented
via the \texttt{bitsandbytes} and \texttt{peft} integration.

1. \textbf{4-bit Quantisation (NF4):} The frozen weights of the pre-trained
model, denoted $W_0$, are compressed into an information-theoretically optimal
non-linear data type: the \textit{NormalFloat 4-bit}. A ``double quantisation''
technique also compresses the normalisation constants, reducing the memory
footprint from 14 GB to approximately 4.5 GB.
2. \textbf{Low-Rank Adaptation (LoRA):} Only additive differential weights
$\Delta W$ are trained, represented as a low-rank matrix decomposition of
rank $r$:
\begin{equation}
    W = W_0 + \Delta W = W_0 + B A
\end{equation}
where $B \in \mathbb{R}^{d \times r}$ and $A \in \mathbb{R}^{r \times k}$
with $r \ll \min(d, k)$. By setting $r=16$ and $\alpha=32$ on the query and
value projection matrices (\texttt{q\_proj, v\_proj}), the number of trainable
parameters dropped to less than $0.5\%$.
3. \textbf{Gradient Checkpointing:} To reduce the memory complexity of the
computational graph (local activation saving), this feature was enabled in the
\texttt{TrainingArguments}, trading VRAM cost for CPU recomputation during
backpropagation.

\subsection{Directed Generation Strategies}
Once the model was made familiar with the ``Original Post $\rightarrow$ Reply''
dynamic, two generation heuristics were tested:

\subsubsection{Analytical Prompt Engineering (\textit{Zero-Shot})}
The first heuristic consists of guiding the model a priori. At initialisation,
the \textit{Prompt Engineering} pipeline analyses the weights $\mathbf{w}$ of
the LinearSVC from Axis 1. By identifying the features with the largest
positive coefficients (e.g., text length or plural pronoun usage), the script
builds a pre-prompt injected dynamically (``Make sure to write a very long and
detailed argument. Use the pronoun `We' extensively.''). Although elegant and
fast (only a single generation is required), this method quickly reaches the
limits of obedience of weakly aligned LLMs.

\subsubsection{Best-of-N Sampling}
The main heuristic deployed is \textit{Search-based generation} or
\textit{Best-of-N}. For a given debate, the stochastic generator produces $N$
disjoint candidate sequences (with probabilistic sampling $T=0.85$ and
$\text{top\_p}=0.9$).
The evaluator then loads the Axis 1 stylistic encoder, computes the Jaccard
distance with the OP and the full profile of each candidate, passes them through
the SVM model, and retains the argument that maximises the decision function.

A \textit{fail-fast} mechanism (Dry-Run) validates the dimensional compatibility
of SVM and encoder tensors (e.g., 384/1024 dim RoBERTa mismatch) in the first
few seconds to safeguard the long inference run. The final result is then
assessed by the RoBERTa judge (Axis 2) to determine its authenticity.

\vspace{0.5cm}
\begin{tcolorbox}[colback=blue!5!white,colframe=blue!75!black,title=Axis 3 Experimental Results]
\textbf{[To be completed: Fill in the average statistics from the generated CSV/HTML report: average Axis 1 score of generated texts, Axis 2 score (Authenticity).]}

\textbf{[To be completed: Include a screenshot of the generated HTML table. Provide 1 or 2 textual examples of particularly convincing generated arguments, and interpret why the Axis 1 SVM judged them relevant (length, vocabulary, structure).]}
\end{tcolorbox}


\section{Project Summary}

\subsection{Challenges Encountered}

The project surfaced a series of technical obstacles that, taken together,
shaped many of our architectural decisions across all three axes.

The most time-consuming hurdle appeared during the fine-tuning of
\textbf{DeBERTa-v3-large} (Axis~1). On top of an architectural incompatibility
with \texttt{AutoModelForMultipleChoice}, resolved by switching to
\texttt{AutoModelForSequenceClassification}, loading the model in
\texttt{float16} caused numerical instability, worked around via
\texttt{bfloat16}. Even after enabling \textit{Gradient Checkpointing} and
reducing the batch size to~1, each epoch still took around twelve minutes on
the RTX~3080. The final accuracy of 60.21\,\% fell below RoBERTa-base's
65.89\,\%, confirming our main takeaway: on a corpus as small as WAC,
the bottleneck is data volume, not model capacity.

Memory pressure was a recurring theme elsewhere. On Axis~2, materialising
TF-IDF output as a dense matrix instantly saturated the 64~GB of RAM on our
compute node; strict use of \texttt{scipy.sparse.csr\_matrix} throughout the
SVM pipeline was a prerequisite for any experiment to run. On Axis~3, loading
Mistral-7B in full precision hit the VRAM ceiling of the RTX~3090 almost
immediately; 4-bit quantisation (QLoRA) combined with \textit{Gradient
Checkpointing} made training feasible on consumer-grade hardware.

Two further issues emerged at the pipeline integration stage. Visualising
generated-text embeddings via PCA produced total overlap; t-SNE revealed
clear clusters but at prohibitive computation times at full scale, resolved
by a \texttt{TruncatedSVD} $\rightarrow$ \texttt{t-SNE} pipeline with
subsampling. More consequentially, saving a trained SVM as a \texttt{.pkl}
without persisting the associated \texttt{TfidfVectorizer} or
\texttt{Word2Vec} vocabulary made the classifiers unusable for inference,
forcing us to fall back on RoBERTa for the Axis~3 evaluation step.

\subsection{Achieved Goals and Gaps}

The three axes (Prediction, Detection, and Generation) were all implemented
and communicate with each other as intended. One unplanned addition was the
\textit{Fail-Fast} (Dry-Run) mechanism introduced in Axis~3, which validates
tensor dimension compatibility on a minimal batch before full-scale GPU
inference; in practice, this reflex saved several hours of wasted computation.

The main deviation from the initial plan concerns the generation strategy.
The original plan included an RLHF-based approach, but its implementation
complexity led the team to adopt \textit{Best-of-N} sampling combined with
\textit{a posteriori} scoring instead, with results convincing enough to
consider the axis validated within a reasonable development time.

\subsection{Lessons Learned and What We Would Do Differently}

Several practices adopted late in the project would have been worth
establishing from day one. The most impactful would have been
\textbf{always saving the complete encoder state} alongside any classifier:
persisting \texttt{TfidfVectorizer} or \texttt{Word2Vec} objects from the
start would have preserved full flexibility in Axis~3. Operationally,
systematic use of \textbf{tmux} or \textbf{nohup} would have prevented SSH
drops from killing long-running scripts on the compute servers.

On the organisational side, clearer ownership of shared pipeline components
from the start, particularly the serialisation conventions for trained
models, would have avoided the inter-axis compatibility issues that only
surfaced at integration time. With five contributors working in parallel,
agreeing on a common I/O contract at the beginning of each axis would have
been more efficient than converging on one under time pressure.

More broadly, this project reinforced how much dataset size constrains what
even a larger model can achieve. The fact that a 400M-parameter model
underperforms a 125M one on 3\,000 training examples is easy to state as a
principle, but becomes concrete when you have spent hours debugging a
pipeline only to confirm it. If we were to redo the project, investing early
in data augmentation or additional collection would be the first priority.

```latex id="j7w2kp"
\section{Team Contributions (Not counted in word limit)}

\textbf{Contributions and Roles:}

\begin{itemize}

    \item \textbf{Wassim SMATI}: I was  responsible for Axis 2 (AI-generated text detection) and Axis 3 (text generation). I designed the project's Hydra-based architecture and configuration management system, ensuring reproducibility and scalability across experiments. For Axis 2, I developed and evaluated AI-generated text detectors on the HC3, GriD, and M4GT datasets, implemented and compared multiple text representations (TF-IDF, Word2Vec, and RoBERTa-based Sentence-Transformers embeddings), and built a RoBERTa fine-tuning pipeline using the Hugging Face Trainer API. For Axis 3, I fine-tuned GPT-2 and Qwen 2.5 3B models using QLoRA under hardware constraints, and implemented both Best-of-N generation and prompt engineering strategies. I also developed interpretability tools, including t-SNE and UMAP visualizations to compare generated and human-written texts.

    \item \textbf{Tristan JIN}: I was mainly responsible for Axis 1: searching for papers and the WAC dataset; proposing the simplification of Axis 1; exploring the dataset and a first failed cleaning attempt. Implementation of the Encoder (W2V, TF-IDF, RoBERTa) + SVM method (tested with LogReg initially) with and without OP; implementation of the paper's features and SHAP analysis; implementation of the pairwise method; failed attempt at RoBERTa fine-tuning.

    \item \textbf{Yanis DAHASSE}: In this project, I focused mainly on Axis 1. I handled the cleaning and formatting of the CMV dataset, as well as data preparation , train/validation/test splitting by post identifier and position bias handling. On the modelling side, I performed the fine-tuning of RoBERTa in Cross-Encoder architecture, which constitutes our best result on this task, and explored several DeBERTa approaches that did not succeed within the allotted time. In parallel, I conducted a literature review of reference works and developed the project timeline.

    \item \textbf{Nicolas CASTEL}:

\end{itemize}
```

\vspace{1cm}

\textbf{Organisation and Tools:}

The work was structured around a shared Git repository. Since Axes 1 and 2 are
largely independent, each member worked on their own branch or on Colab, without
requiring frequent synchronisation. Code merging onto the \texttt{main} branch
took place during Axis 3, which requires the outputs of the two previous axes.
Hyperparameter management was handled by the \textbf{Hydra} framework, allowing
configurations to be shared via \texttt{YAML} files without introducing conflicts
in the Python code. Heavy model execution was distributed across Télécom Paris
GPUs, Google Colab and Onyxia depending on needs and availability.




% =========================================
% SECTION: AI TOOLS USAGE
% ==========================================
\section{Use of AI Tools (Not counted in word limit)}

We indicate here how we used AI tools
(Claude, Gemini, GitHub Copilot) during the project.

\begin{itemize}
    \item \textbf{In the code:} AI tools were used for debugging
    and object-oriented code refactoring, as well as for writing repetitive
    functions such as stylistic feature extraction. All code was nevertheless
    reviewed, tested and manually executed.
    \item \textbf{In the report:} AI tools facilitated LaTeX formatting
    (tables, formulas, figures) as well as spell-checking and grammar
    correction. They also helped with the presentation of complex models such
    as RoBERTa. All analyses, interpretations and discussions of results
    remain the product of our own thinking.
\end{itemize}
% ==========================================
% BIBLIOGRAPHY
% ==========================================
\newpage
\begin{thebibliography}{9}

\bibitem{tan2016winning}
C. Tan, V. Niculae, C. Danescu-Niculescu-Mizil, and L. Lee,
\textit{Winning Arguments: Interaction Dynamics and Persuasion Strategies in Good-faith Online Discussions},
Proceedings of the 25th International Conference on World Wide Web (WWW 2016).

\bibitem{hc3}
B. Guo et al.,
\textit{How Close is ChatGPT to Human Experts? Comparison Corpus, Evaluation, and Detection},
arXiv preprint arXiv:2301.07597 (2023).

\bibitem{qlora}
T. Dettmers, A. Pagnoni, A. Holtzman, and L. Zettlemoyer,
\textit{QLoRA: Efficient Finetuning of Quantized LLMs},
Advances in Neural Information Processing Systems (NeurIPS 2023).

\bibitem{attention}
A. Vaswani et al.,
\textit{Attention Is All You Need},
Advances in Neural Information Processing Systems (NeurIPS 2017).

\bibitem{tsne}
L. van der Maaten and G. Hinton,
\textit{Visualizing Data using t-SNE},
Journal of Machine Learning Research (JMLR 2008).

\end{thebibliography}

\end{document}